Adversarial prompt evaluation currently sits at the intersection of jailbreak datasets, multilingual safety corpora, and broader dataset-governance practice. Public resources such as All-Prompt-Jailbreak, WildJailbreak, Aya RedTeaming, and XSafety are useful starting points because they make attack prompts easier to compare across systems, but they also illustrate the fragmentation of the space: some are template-heavy, some are multilingual without deep native adversarial coverage, and some are broader safety corpora rather than prompt-classifier benchmarks \citep{allpromptjailbreak2024,wildjailbreak2024,ayaredteaming2024,xsafety2024}.

From a benchmark-construction perspective, \system adopts the view that the artifact should be treated as a governed dataset bundle rather than as an informal collection of CSV rows. This aligns with the dataset documentation and reproducibility discipline emphasized by datasheet-style governance work \citep{gebru2021datasheets}. The additional requirement here is that reproducibility must hold across transformations, multilingual expansion, and evaluation-specific perturbations rather than only over a static released file.

The main gap this paper targets is not simply ``more jailbreak data.'' It is the lack of a benchmark grammar that combines multilingual coverage, leakage-resistant generation, and explicit accounting for utility and efficiency costs. In that sense, \system is closer to an evaluation instrument than to a one-shot corpus release \citep{shibbolethdraft2026}.
